\subsection{A compactification system independent of the chosen model}

The definition on 049(c) can retain every smooth normal-crossings model
at once. This gives a choice-independent object even when no individual
model is distinguished by the open variety.

\begin{theorem}[Proved: the canonical system of smooth compactifications]
\label{A6-model-system}
Let $k$ be a field of characteristic zero and let $U/k$ be smooth,
geometrically connected, separated and of finite type. Define
$\mathcal C_k(U)$ as follows. An object is a smooth proper integral
$k$-scheme $X$, with a specified dense open immersion $j_X:U\hookrightarrow X$
whose reduced complement $D_X$ is a strict normal-crossings divisor.
A morphism from $X$ to $Y$ is a $k$-morphism $f:X\to Y$ satisfying
$f\circ j_X=j_Y$. Then:
\begin{enumerate}[label=(\roman*)]
\item $\mathcal C_k(U)$ is essentially small, nonempty and cofiltered.
There is at most one morphism between any two objects. Every morphism
is proper and satisfies $f^{-1}(j_Y(U))=j_X(U)$.
\item Its inverse system, with its marked common interior,
\[
 \mathfrak C_k(U)=\bigl((X,D_X)\bigr)_{X\in\mathcal C_k(U)},
 \qquad U\longrightarrow\mathfrak C_k(U),
\]
is independent, up to canonical pro-isomorphism, of the choice of a
small skeleton. For any fixed model $X_0$, restricting to all models
dominating $X_0$ gives the same pro-object. The same holds after fixing
any finite collection of initial models.
\item Every $k$-isomorphism $\varphi:U\to V$ extends uniquely to an
isomorphism of these pro-systems with their marked interiors. Thus,
on any class $\mathcal A$ for which the augmented arithmetic fundamental
group is fully faithful on isomorphisms, that group reconstructs
$\mathfrak C_k(U)$ functorially. The reconstructed object retains the
schemes, boundaries and transition maps, not just the abstract indexing
partially ordered set.
\end{enumerate}
Here an isomorphism of marked pro-systems includes the isomorphism of
their interiors and commutes with the displayed maps from those
interiors. The theorem makes no selection of one compactification.
\end{theorem}
\begin{proof}
Since $U$ is smooth and geometrically connected, it is geometrically
integral. Nonemptiness is Theorem~\ref{A6-SNC-existence}. Essential
smallness follows because finite-type $k$-schemes, finite affine gluing
data and morphisms between them form a set up to isomorphism; take any
small skeleton when forming the inverse system.

Two morphisms $X\to Y$ agreeing on $U$ are equal: the equalizer is closed
because $Y$ is separated, and equality on the dense open determines a
morphism from the reduced scheme $X$. Every such morphism is proper.
Indeed its graph is closed in $X\times_kY$, whose projection to $Y$ is
proper. Put $W=f^{-1}(j_Y(U))$. The given copy of $U$ in $W$ is a section
of the separated morphism $W\to U$, hence is closed in $W$; it is also
open. The nonempty open $W$ of the integral scheme $X$ is irreducible,
so this open-and-closed copy of $U$ is all of $W$. This proves the asserted
preimage equality. In particular the divisorial log structures of the
pairs are preserved by all transition maps.

For two models $X_1,X_2$, let $\Gamma$ be the reduced closure of the
diagonal copy of $U$ in $X_1\times_kX_2$. It is integral and proper over
$k$. The intersection with $j_{X_1}(U)\times_kX_2$ is exactly the graph
of $j_{X_2}$, because that graph is closed over $U$ by separatedness.
Consequently $U$ is open in $\Gamma$. Let $Z=\Gamma\setminus U$ with
its reduced closed subscheme structure. Every singular point of
$\Gamma$ lies in $Z$. Apply Temkin
\cite[Theorem~1.1.11, p.~4]{Temkin2018A567} to $(\Gamma,Z)$.
Its hypotheses hold: $\Gamma$ is noetherian and quasi-excellent, being
of finite type over the characteristic-zero field $k$. The theorem
gives a modification $X_3\to\Gamma$ supported on $Z$, with $X_3$
regular and with total transform of $Z$ strictly monomial. Its reduced
support is therefore a strict normal-crossings divisor. The modification
is an isomorphism over $U$; $X_3$ is proper over $k$, and regularity over
the perfect field $k$ gives smoothness. The two projections make $X_3$
an object dominating $X_1$ and $X_2$ in $\mathcal C_k(U)$.

This proves the common-domination condition. Parallel arrows are already
equal, so the equalizing condition for a cofiltered category holds as
well. For a fixed $X_0$, every model $X$ has a common refinement with
$X_0$, and any two such refinements have a further common refinement.
Thus the models over $X_0$ are cofinal in the precise inverse-system
sense: the category of their arrows to any given $X$ is nonempty and
connected. Iterating gives the assertion for finitely many fixed models.
The defining formula
\[
 \operatorname{Hom}_{\mathrm{Pro}}
       ((X_i),(Y_j))
   =\varprojlim_j\varinjlim_i\operatorname{Hom}(X_i,Y_j)
\]
then proves that these restrictions and changes of skeleton give
canonical pro-isomorphisms.

Let $\varphi:U\to V$ be an isomorphism. For each model $Y$ of $V$,
regard $Y$ as a model of $U$ using the marking $j_Y\circ\varphi$.
The identity map of that scheme gives a representative of a pro-map
$\mathfrak C_k(U)\to\mathfrak C_k(V)$. These representatives are
compatible with every transition map. Equivalently, to use any prescribed
initial source model, resolve the closure of the graph of $\varphi$ in
its product with $Y$, by the common-domination construction above.
Two representatives extending $\varphi$ agree on a common refinement:
they agree on its dense open $U$ and the target is separated. Hence the
pro-map is unique. Applying the construction to $\varphi^{-1}$ and using
this uniqueness makes both composites the identity. It also proves
compatibility with composition.

For $U,V\in\mathcal A$, a compatible outer isomorphism of augmented
arithmetic groups first yields its unique isomorphism $\varphi:U\to V$
by the defining Isom-anabelian hypothesis. The construction just proved
extends that particular isomorphism uniquely. Conversely an isomorphism
of the marked pro-systems restricts to its isomorphism of interiors,
which induces the original group isomorphism. This is a deduction from
the stated anabelian input; it imposes no anabelian hypothesis on the
individual boundary strata. The strata, their normal bundles and the
normal-torus families of tangential cover functors are reconstructed
for each model because they are intrinsic constructions from that
scheme and its boundary. This recovers the family of possible frames,
not a distinguished frame or point of a normal torus.
\end{proof}

The supported modification used here is provided by Temkin's theorem.
For a description by successive local operations, the birational map
between any two smooth proper models also has the weak factorization
of Abramovich--Karu--Matsuki--W{\l}odarczyk
\cite[Theorem~0.3.1, especially~(6), pp.~2--3]{AKMW2002A567}:
the base field may be any characteristic-zero field, the varieties are
complete and nonsingular, and the centers lie outside $U$ and have
normal crossings with the relevant boundary transforms. This is a
zigzag of blow-ups and blow-downs. The common-domination proof above
requires no assertion that both maps from its common model are
composites of blow-ups along smooth centers.

\begin{theorem}[Proved: compatibility of tangential specialization across models]
\label{A6-model-specialization}
Let $k$ have characteristic zero, put $R=k[[t]]$ and $F=k((t))$, and fix
a smooth geometrically connected separated finite-type $F$-scheme $U$.
Consider proper integral regular $R$-schemes $\mathcal X$ containing
$U$ as a dense open, with reduced complement $D_{\mathcal X}$ a strict
normal-crossings divisor. Give them their divisorial log structures,
and give $S=\operatorname{Spec}R$ the log structure of $t=0$.
Their morphisms are $R$-morphisms restricting to $\operatorname{id}_U$.
This category is nonempty and cofiltered. Every such model is flat
over $R$, and it has a canonical morphism of log schemes to $S^{\log}$.

For every model, write $\mathcal X_0^{\log}$ for its strict closed
log fibre. The equivalences
\[
 \operatorname{F\acute Et}(U)
 \simeq\operatorname{K\acute Et}(\mathcal X^{\log})
 \simeq\operatorname{K\acute Et}(\mathcal X_0^{\log})
\]
are compatible with every transition map. In particular, if
$f:\mathcal X'\to\mathcal X$ is one, then
\[
 f_0^*:\operatorname{K\acute Et}(\mathcal X_0^{\log})
       \longrightarrow
       \operatorname{K\acute Et}(\mathcal X_0'{}^{\log})
\]
is an equivalence. The comparisons between any two models are independent
of common refinement, up to the unique natural isomorphism compatible
with restriction to $U$. They preserve the maps to $G_F$ and, for the
same chosen frame $t$, the tangential specialization extensions over
$G_k$ of Theorem~\ref{A6-log-specialization}.
\end{theorem}
\begin{proof}
The ring $R$ is excellent and noetherian. The map $U\to S$ is separated
and of finite type, since $F=R[1/t]$ is of finite type over $R$.
Apply Temkin \cite[Theorem~1.1.15, p.~5]{Temkin2018A567} with
$Y_0=Y=U$ and with base $S$. It produces a regular proper model with
the required strict normal-crossings complement. Taking the component
containing $U$ removes any extraneous components; in the construction
in \cite[\S3.5.3, p.~24]{Temkin2018A567}, $U$ is already dense in its
compactification. This proves nonemptiness. All the schemes in the
stated category are integral and dominate $S$. Their coordinate rings
are torsion-free $R$-modules and hence are flat over the discrete
valuation ring $R$. The function $t$ has divisor supported on the
boundary. Locally it is a unit times a product of powers of boundary
parameters, so it defines the indicated morphism of divisorial log
schemes over $S^{\log}$.

The common-domination proof of Theorem~\ref{A6-model-system} works over
$R$: take the closure of the graph of $U$ in
$\mathcal X_1\times_R\mathcal X_2$ and apply
\cite[Theorem~1.1.11]{Temkin2018A567} to its closed complement of $U$.
This graph closure is a noetherian quasi-excellent characteristic-zero
scheme, because it is of finite type over the excellent ring $R$.
Its resolution is regular, proper and integral over $R$, is unchanged
on $U$, and has strict normal-crossings boundary. It is again flat by
the torsion-free argument. Uniqueness of arrows and the equality
$f^{-1}(U)=U$ follow from the same density and section argument as
before. Thus all arrows induce canonical log morphisms over $S^{\log}$.

A regular scheme with its strict normal-crossings divisorial log
structure is log regular. The cover equivalences are therefore exactly
the application of Hoshi's
\cite[Proposition~A.10 and Corollary~3.9]{HoshiLog2009A567}
proved in Theorem~\ref{A6-log-specialization}: every residue
characteristic is zero, so boundary ramification of a finite cover of
$U$ is tame, and properness over the complete local base supplies the
closed-fibre equivalence. All these hypotheses hold for every model
in the present category; semistability or a saturated log-smooth
structure is not being silently assumed.

For an arrow $f$, restriction and pullback give the commuting diagram
\[
 \begin{array}{ccc}
 \operatorname{K\acute Et}(\mathcal X^{\log})&
   \xrightarrow{\ f^*\ }&
 \operatorname{K\acute Et}(\mathcal X'{}^{\log})\\
 \downarrow&&\downarrow\\
 \operatorname{F\acute Et}(U)&\xrightarrow{\ \operatorname{id}\ }&
 \operatorname{F\acute Et}(U).
 \end{array}
\]
The vertical functors are equivalences, so $f^*$ is one. The analogous
square of restrictions to the strict closed log fibres then shows that
$f_0^*$ is an equivalence. More precisely, if $j_{\mathcal X}$ and
$i_{\mathcal X}$ denote the interior and closed-fibre inclusions, equip
each special-fibre cover category with the comparison
\[
 E_{\mathcal X}=j_{\mathcal X}^*\circ(i_{\mathcal X}^*)^{-1}
       :\operatorname{K\acute Et}(\mathcal X_0^{\log})
        \longrightarrow\operatorname{F\acute Et}(U),
\]
including the natural isomorphisms of the chosen quasi-inverse.
Every comparison between models is equipped with its isomorphism to
the identity under these $E_{\mathcal X}$. Full faithfulness supplies
a unique natural isomorphism between any two comparisons that respects
these specified identifications. Bare uniqueness without that
normalization is not asserted. Consequently the comparisons
are independent of refinement, and their isomorphisms obey the cocycle
condition. This is the stated categorical meaning of canonicity.

All functors in the two squares commute with pullback from $S^{\log}$.
Thus their induced outer fundamental-group isomorphisms commute with
the augmentations to $G_F$. Keeping the same section
$\sigma_t:G_k\to G_F$ and pulling back these augmented groups gives
the same tangential extension over $G_k$. Changing the frame has the
Kummer dependence already computed in Theorem~\ref{A6-log-specialization};
changing the model introduces no additional ambiguity. The construction
is for models of this fixed generic open over this fixed trait. It
does not assert that an unrelated compactification over $k$ is a model
of that family.
\end{proof}

\begin{proposition}[Refuted: reconstruction of a chosen special fibre]
\label{A6-special-fibre-counterexample}
Over $R=\Q[[t]]$, there are two proper integral regular models of the
same smooth geometrically connected generic open, both with strict
normal-crossings reduced boundary, whose special fibres are not
isomorphic. Nevertheless their strict special log fibres have equivalent
finite Kummer-\'etale cover categories over $G_{\Q((t))}$ and the same
tangential arithmetic extension over $G_\Q$ for the fixed frame $t$.
Thus even these augmented and framed cover data do not reconstruct a
chosen special-fibre scheme. This does not contradict reconstruction of
the complete marked model system.
\end{proposition}
\begin{proof}
Put $F=\Q((t))$, take
\[
 \mathcal X=\mathbf P^1_R\times_R\mathbf P^1_R,
 \qquad
 U=(\mathbf P^1_F\setminus\{0,1,\infty\})^2,
\]
and let $D$ be the union of the six horizontal boundary divisors and
the vertical fibre $\mathcal X_0$. This is a strict normal-crossings
divisor on the proper regular integral $R$-scheme $\mathcal X$.
Blow up the closed rational point $P=(x,y,t)=(0,0,0)$ and write
\[
 f:\mathcal X' =\operatorname{Bl}_P\mathcal X\longrightarrow\mathcal X,
 \qquad D'=(f^{-1}D)_{\mathrm{red}}.
\]
The blow-up is projective, is an isomorphism away from $P$, and fixes
$U$. Near $P$ the three boundary parameters are $x,y,t$. Its charts are
\[
 \begin{array}{c|c|c|c}
 \text{chart}&\text{substitution}&D'&t\text{ on the chart}\\\hline
 x&y=xv,\ t=xu&xvu=0&xu\\
 y&x=ya,\ t=yb&yab=0&yb\\
 t&x=tc,\ y=td&tcd=0&t .
 \end{array}
\]
The blow-up of the regular parameter ideal $(x,y,t)$ is regular:
the completed local chart equations $t-xu$, respectively $t-yb$,
eliminate $t$ and leave regular power-series rings; the $t$-chart is
regular directly. Away from the exceptional divisor the map is
an isomorphism. The vanishing factors in the table form part of a
regular parameter system, proving the strict normal-crossings
condition. The other horizontal divisors miss $P$ and are unchanged
near the exceptional locus. The total space $\mathcal X'$ is integral
and dominates $\operatorname{Spec}R$, so it is flat over this discrete
valuation ring. Both models therefore satisfy every hypothesis of
Theorem~\ref{A6-model-specialization}.

The old special fibre is the irreducible surface
$\mathcal X_0=\mathbf P^1_\Q\times\mathbf P^1_\Q$.
Since $t$ vanishes to order one at $P$, the new special fibre is reduced
and is the union
\[
 \mathcal X'_0=
 \operatorname{Bl}_{(0,0)}
       (\mathbf P^1_\Q\times\mathbf P^1_\Q)
 \ \cup_{\mathbf P^1_\Q}\ \mathbf P^2_\Q .
\]
The second component is the exceptional divisor
$\mathbf P(T_P\mathcal X)$; its intersection with the first is
$\mathbf P(T_P\mathcal X_0)$. The two factors $xu$, respectively $yb$,
in the table verify multiplicity one for both components. Thus the
special fibres have respectively one and two geometric irreducible
components and cannot be isomorphic, even after extending $\Q$.

Give both total spaces their boundary log structures and both special
fibres their strict induced log structures. The cited theorem identifies
their Kummer-\'etale cover categories by $f_0^*$, with the identifications
normalized by their common restriction to $\operatorname{F\acute Et}(U)$.
These identifications commute with pullback from the standard log trait.
Consequently their fundamental groups are isomorphic over the same
$G_F$, up to the usual basepoint conjugacy. Pulling back along the same
root-fixing section $\sigma_t:G_\Q\to G_F$ gives identical augmented
extensions with geometric kernel $\pi_1(U_{\overline F})$ over $G_\Q$.
In particular the associated tangential outer actions agree, although
the special-fibre schemes differ. The algebraic model and blow-up also
descend to $\Q[t]_{(t)}$; the generic open over $\Q(t)$ is a
dimension-two hyperbolic polycurve. No unproved anabelianity hypothesis
is needed for this counterexample or for its cover comparison.
\end{proof}
